% !TeX root = ../../../main.tex
\section{Weighted Witt addition in the cyclic calculation}

The following homogeneity property of the universal addition polynomials is what
makes degree estimates on Witt sums possible: even though the carry polynomials $C_n$
of \Cref{eq:witt_addition_shape} are complicated, they are \emph{isobaric} for the
weighting in which the $j$-th component has weight $p^j$.

\begin{lemma}[isobaricity of Witt addition]\label{lemma:witt_addition_isobaric}
Assign to the variables $X_j$ and $Y_j$ the weight $p^j$. Then the $n$-th addition
polynomial $S_n$ is isobaric of weight $p^n$: every monomial of $S_n$ has total weight
exactly $p^n$. More generally, write
$\vec a_i=(a_{i,0},\dots,a_{i,r-1})$ and assign $a_{i,j}$ weight $p^j$.
Then, for every $0\leq n<r$, the $n$-th component of the iterated Witt sum
$\bigoplus_{i=1}^{d}\vec a_i$ is an isobaric polynomial of weight $p^n$ in the
variables $a_{i,j}$.
\end{lemma}

\begin{proof}
We first prove the assertion for two summands by induction on $n$.
For $n=0$, we have $S_0=X_0+Y_0$, which is isobaric of weight $1$.
Suppose now that $S_j$ is isobaric of weight $p^j$ for every $j<n$.
The ghost identities give
\[
p^n S_n = \sum_{j \le n} p^j \bigl(X_j^{p^{n-j}} + Y_j^{p^{n-j}}\bigr)
- \sum_{j < n} p^j S_j^{p^{n-j}}
\]
Each $X_j^{p^{n-j}}$ and $Y_j^{p^{n-j}}$ has weight
$p^j p^{n-j}=p^n$. Also, every monomial in the expansion of
$S_j^{p^{n-j}}$ is a product of $p^{n-j}$ monomials of weight $p^j$,
so it has weight $p^n$. Thus the right-hand side is isobaric of weight
$p^n$. This identity holds in the universal polynomial ring over $\Z$,
where multiplication by the nonzero integer $p^n$ does not change which
monomials occur. Hence $S_n$ is isobaric of weight $p^n$ as well.

We next induct on the number $d$ of summands. Write
\[
\vec A^{(d)}:=\bigoplus_{i=1}^{d}\vec a_i
=\vec a_1\oplus\cdots\oplus\vec a_d,
\]
where each variable $a_{i,j}$ has weight $p^j$. For $d=1$, the component
$A_n^{(1)}=a_{1,n}$ has weight $p^n$. Suppose the assertion holds for
$d-1$ summands. Then
\[
A_n^{(d)}=
S_n\bigl(A_0^{(d-1)},\dots,A_n^{(d-1)};
a_{d,0},\dots,a_{d,n}\bigr).
\]
By induction, $A_j^{(d-1)}$ is isobaric of weight $p^j$.
Thus each substituted polynomial has the same weight as the variable it
replaces, so substitution preserves weights. Since $S_n$ is isobaric of
weight $p^n$, so is $A_n^{(d)}$, completing the induction.
\end{proof}
